\chapter{栈（Stack）：后进先出的智慧}

在学习了通用的“线性表”之后，我们进入一类特殊的线性表——\emph{受限的线性表}。所谓受限，并不是功能变弱了，而是通过规则的约束，让它在特定场景下变得极其强大。栈就是这样一种结构。

\section{什么是栈？——后进先出的世界}

栈（Stack）是一个特殊的线性表，它只允许在\emph{一端}进行插入和删除操作。这一端称为\emph{栈顶}，另一端称为\emph{栈底}。不含任何元素的栈称为\emph{空栈}。

\subsection{核心原则：后进先出（LIFO）}

栈的精髓是 \emph{LIFO（Last In First Out）}，即后进先出。最后放入栈中的元素，最先被取出。

\textbf{生活中的类比}：

\begin{itemize}
    \item \emph{叠盘子}：你洗好一个盘子就叠在最上面；要用盘子时也是从最上面拿。最后放上去的盘子最先被拿走。
    \item \emph{步枪弹夹}：子弹一颗一颗压进去（入栈），射击时最上面那颗最先出膛（出栈）。
    \item \emph{浏览器后退按钮}：你访问的页面按顺序被压入一个栈，点击“后退”时，最后访问的页面最先被退出。
\end{itemize}

\subsection{为什么是“受限”的？}

与普通线性表相比，栈的操作被严格限制在一端。这种限制带来了两个好处：

\begin{itemize}
    \item \emph{逻辑清晰}：我们只关心栈顶元素，无需考虑中间元素的随机访问。
    \item \emph{实现简单}：只需要维护一个栈顶指针，插入和删除操作都集中在同一端。
\end{itemize}

\section{栈的基本操作}

栈的操作非常简洁，主要就是两个动作：

\begin{itemize}
    \item \emph{入栈（Push）}：将一个元素放到栈顶，成为新的栈顶。
    \item \emph{出栈（Pop）}：取出当前栈顶元素，栈顶指针下移，原来的次顶元素成为新栈顶。
\end{itemize}

\subsection{Python 中的栈实现}

Python 内置的 \texttt{list} 类型就是一个天然的栈：\texttt{append()} 对应入栈，\texttt{pop()} 对应出栈。我们可以用它来模拟浏览器的前进后退：

\begin{minted}{python}
# 初始化一个空栈，记录浏览历史
history_stack = []

# 模拟访问页面（入栈）
history_stack.append("首页")
history_stack.append("搜索页")
history_stack.append("详情页")
print(f"当前浏览路径: {history_stack}")

# 模拟点击“返回”按钮（出栈）
if history_stack:
    last_page = history_stack.pop()
    print(f"离开: {last_page}")
    print(f"回到上一页: {history_stack[-1] if history_stack else '空'}")
\end{minted}

运行这段代码，你会看到“详情页”被弹出，当前页面回到了“搜索页”。这正是后进先出的体现。

\section{栈的两种存储结构}

根据底层存储方式的不同，栈可以分为顺序栈和链栈。

\subsection{顺序栈}

顺序栈使用一组连续的存储单元（如数组）来存放元素，栈顶通常用一个整型变量 \texttt{top} 来指示。

\begin{itemize}
    \item \emph{优点}：访问速度快，实现简单。
    \item \emph{缺点}：需要预先分配空间，当栈满时无法再插入新元素，会发生\emph{栈溢出（Stack Overflow）}。
\end{itemize}

\begin{quote}
\emph{小知识}：“栈溢出”这个词大家可能经常在编程错误中见到，它就是因为递归调用太深或局部数据太大，导致系统分配的栈空间不够用了。
\end{quote}

\subsection{链栈}

链栈使用链表结构，每个节点包含数据和指向下一个节点的指针。栈顶就是链表的头节点。

\begin{itemize}
    \item \emph{优点}：动态分配内存，理论上不会满（除非内存耗尽）。
    \item \emph{缺点}：每个节点需要额外的指针空间，访问速度略慢于顺序栈。
\end{itemize}

现代高级语言（如 Java、Python）的虚拟机在管理方法调用时，底层大量使用了链式栈的思想，以便灵活应对复杂的调用层次。

\section{栈的经典应用}

栈不仅仅是一种存储数据的方式，它代表了一种\emph{“保存现场、稍后回来”}的思维模式。下面我们看看栈在现实和算法中的几个典型应用。

\subsection{撤销操作（Undo）}

在文字处理软件（如 Word）中，你的每一次编辑（输入文字、删除、格式化）都会被压入一个“撤销栈”。当你按下 Ctrl+Z 时，系统从栈中弹出最近的操作，并执行反向动作，从而恢复到之前的状态。这正是后进先出思想的体现——最后做的操作最先被撤销。

\subsection{括号匹配}

编译器如何判断代码中的括号是否成对匹配？比如 \texttt{(( ))} 是合法的，而 \texttt{({ )}} 是不合法的。算法很简单：

\begin{itemize}
    \item 从左到右扫描字符串，遇到左括号（\texttt{(}、\texttt{\{}、\texttt{[}）就将其\emph{入栈}。
    \item 遇到右括号（\texttt{)}、\texttt{\}}、\texttt{]}）时，从栈顶\emph{弹出一个左括号}，检查是否匹配。
    \item 如果栈为空却遇到右括号，或者最后栈中还有剩余左括号，说明匹配失败。
\end{itemize}

这个算法只需要一次遍历，时间复杂度 O(n)，空间复杂度 O(n)（最坏情况全是左括号）。很多编程面试题都喜欢考察这个题目。

\subsection{表达式求值（逆波兰表达式）}

计算器在处理 \texttt{1 + 2 * 3} 时，并不是简单地从左到右计算，因为乘法的优先级高于加法。栈可以帮助我们暂存操作数和运算符，从而正确处理优先级。

一种经典方法是先将中缀表达式转换为后缀表达式（逆波兰表达式），然后用一个栈来求值。例如，\texttt{1 + 2 * 3} 对应的后缀表达式是 \texttt{1 2 3 * +}，求值过程如下：

\begin{itemize}
    \item 遇到操作数 \texttt{1}、\texttt{2}、\texttt{3} 依次入栈。
    \item 遇到运算符 \texttt{*}，弹出两个操作数 \texttt{2} 和 \texttt{3}，计算 \texttt{2 * 3 = 6}，将 \texttt{6} 入栈。
    \item 遇到运算符 \texttt{+}，弹出 \texttt{6} 和 \texttt{1}，计算 \texttt{1 + 6 = 7}，得到结果。
\end{itemize}

整个过程只需要一个栈，非常简洁。

\subsection{函数调用与递归}

在程序运行时，每当调用一个函数，系统就会在内存的“调用栈”上为该函数分配一块区域（称为栈帧），用来存放局部变量、参数和返回地址。当函数执行完毕，对应的栈帧被弹出，控制权交还给调用者。

递归函数正是利用了栈的这种特性：每次递归调用都会压入一个新的栈帧，返回时再弹出。如果递归深度过大，就会导致栈溢出。

\section{深度探索：迷宫寻路与回溯}

栈在人工智能和路径规划中的一个经典应用是\emph{迷宫寻路}。假设我们在一个二维迷宫中，从起点出发，要找到一条通往终点的路径。我们可以用栈来记录走过的路径，并在遇到死胡同时回溯。

\textbf{基本思路}：

\begin{enumerate}
    \item 从起点出发，将当前坐标压入栈。
    \item 探测四个方向（上、下、左、右）：
    \begin{itemize}
        \item 如果某个方向可以走（未走过且不是墙），就将新坐标压栈，并标记为已走。
        \item 如果所有方向都不可走，说明当前是死胡同，则\emph{出栈}（回溯）到上一个位置，尝试其他方向。
    \end{itemize}
    \item 重复步骤 2，直到到达终点或栈空（无路可走）。
\end{enumerate}

这种“试探-回溯”的过程，正是\emph{深度优先搜索（DFS）}的核心思想，而栈则是实现这一思想的天然工具。

\begin{quote}
为什么栈适合回溯？\\
因为回溯需要回到最近访问的点，而栈的后进先出特性恰好能记住最近的选择。每一步的决策都被压入栈，当发现此路不通时，只需弹出栈顶，就能回到上一个分叉点。
\end{quote}

\section{栈与队列的对比}

在数据结构中，栈和队列经常被放在一起讨论，因为它们都是受限的线性表，但规则恰好相反：

\begin{table}[h]
\centering
\begin{tabular}{|l|l|l|}
\hline
\textbf{特性} & \textbf{栈（Stack）} & \textbf{队列（Queue）} \\
\hline
原则 & 后进先出（LIFO） & 先进先出（FIFO） \\
操作端 & 同一端（栈顶） & 两端（队尾插入，队头删除） \\
适用场景 & 回溯、撤销、递归 & 排队、任务调度、广度优先搜索 \\
\hline
\end{tabular}
\caption{栈与队列的对比}
\end{table}

简单来说，栈是为了“记住过去”（回溯），而队列是为了“保证公平”（先进先出）。

\section{自测练习}

为了巩固对栈的理解，这里给出几道练习题。你可以尝试用笔算，也可以动手写代码验证。

\begin{enumerate}
    \item \textbf{出栈序列判断} \\
    若入栈序列为 \texttt{A, B, C}，且允许在任意时刻出栈，下列哪个序列是不可能出现的出栈序列？ \\
    A. \texttt{C, B, A} \\
    B. \texttt{B, A, C} \\
    C. \texttt{C, A, B} \\
    D. \texttt{A, B, C}

    \item \textbf{括号匹配加强版} \\
    给定一个只包含 \texttt{(}、\texttt{)}、\texttt{\{}、\texttt{\}}、\texttt{[}、\texttt{]} 的字符串，判断它是否有效。要求使用栈实现。

    \item \textbf{逆波兰表达式求值} \\
    输入一个逆波兰表达式（后缀表达式），例如 \texttt{["2","1","+","3","*"]}，计算其结果。

    \item \textbf{模拟浏览器的前进后退} \\
    设计一个浏览器历史记录类，支持访问新页面、后退、前进操作。用两个栈来实现。

    \item \textbf{最小栈} \\
    设计一个栈，除了支持常规的 push、pop 操作外，还要能在常数时间内获取栈中的最小元素。

    \item \textbf{用栈实现队列} \\
    使用两个栈来实现一个队列，支持入队和出队操作。

    \item \textbf{回文链表检测} \\
    给定一个单链表，判断它是否回文（例如 \texttt{1->2->2->1}）。可以借助栈实现：先遍历一遍将值入栈，再第二次遍历时与栈顶比较。

    \item \textbf{迷宫问题} \\
    用栈实现一个简单的迷宫求解程序，输出从起点到终点的路径。
\end{enumerate}

\section*{本章小结}

栈是一种“简约而不简单”的数据结构。它在计算机底层（函数调用）、算法（深度优先搜索）和日常应用（撤销、括号匹配）中无处不在。掌握了栈，你就掌握了程序如何“记住过去”的方法。

在下一章中，我们将学习另一种受限的线性表——\emph{队列}，看看它又是如何实现“公平”的。
